\section*{产业研发经历}

\noindent
\textbf{中国电信研究院网络研发中心} \hfill 中国北京 / 广州\\
\textit{应用研发工程师}
   \hfill 2021年--2022年 \\
- 构建基于 LSTM 的交通流预测模型，利用实时传感器数据进行短时交通流预测。\\
- 开发基于强化学习的信号配时优化模型，用于提升路口交通信号灯控制效果。\\
- 基于自有数据训练车牌检测与识别模型，使用 YOLO 完成检测，并结合 OCR 模型读取车牌内容。\\
- 参与 MQTT 数据管道研发，覆盖发布/订阅与连接处理，为模型接入实时边缘与路侧数据；同时参与前端产品设计与需求拆解，并与产品和平台团队协作。\\

\noindent
\textbf{中国电信研究院云网研究所} \hfill 中国广州\\
\textit{云网研究员}
   \hfill 2020年--2021年\\
- 起草面向云网架构的 ITU-T 测试与监测框架初稿，将需求转化为测试方法、测试用例、指标体系与标准化报告格式，并通过内部评审与修订持续完善。\\
- 构建基于 Prometheus 的原型系统，用于框架实践运行，支持批量测试执行与关键指标的持续采集。\\

\noindent
\textbf{广发证券} \hfill 中国广州\\
\textit{量化研究员}
   \hfill 2018年\\
- 面向基金管理开展低风险资产配置研究，使用回归分析与线性优化方法，将风险与合规约束转化为可执行策略。\\
- 构建回测与风险控制系统原型，覆盖数据处理、机器学习因子提取（包括 LSTM）与策略执行，并用于历史数据上的策略测试。
